% Section 2 — Preliminaries (charter §4.2; issue T1-1, WOV-262)
\section{Preliminaries: a measure-theoretic setup}
\label{sec:prelim}

% Area role (charter §2): rigorous setup so that the later DPI / sufficiency
% arguments (T2-2) are exact rather than heuristic. Everything downstream --
% information parity (§4), the bounded-learner gap (§5), the parity defect
% (§6) -- is phrased against the objects fixed here. Symbols live in
% notation.tex; this section only instantiates them.

This section fixes the measure-theoretic objects the rest of the paper
quantifies over: the data-generating measure, a \emph{representation} as a
measurable map, a \emph{predictor} as a Markov kernel, and the loss/risk
functionals. The aim is modest but load-bearing: once representations and
predictors are measurable objects on a common probability space, the
data-processing and sufficiency arguments of \cref{sec:parity} become exact
statements about conditional expectations rather than informal claims about
``information.'' We work throughout on standard Borel spaces, which is what
licences the disintegration machinery in \cref{subsec:disint}; the assumption
is harmless for our objects (finite token alphabets, countable program
syntax, and their Polish completions). We follow the standard treatments of
\citet{folland1999real}, \citet{kallenberg2002foundations}, and
\citet{cinlar2011probability}.

\subsection{The data-generating measure}
\label{subsec:dgm}

Fix a probability space $(\Omega, \mathcal{F}, \Prob)$ on which all random
elements are defined. The \emph{input space} $\Xspace$ (a program, or a
prompt-rendered code artifact) and the \emph{target space} $\Yspace$ (the
quantity to be predicted: a label, a program property, a continuation) are
each equipped with a $\sigma$-algebra; we take both to be standard Borel
spaces and write $\Bor{\Xspace}$ and $\Bor{\Yspace}$ for their Borel
$\sigma$-algebras.

\begin{assumption}[Standard Borel data model]
\label{ass:sbs}
$\Xspace$ and $\Yspace$ are standard Borel spaces. The data are a random
element $(X, Y) : \Omega \to \Xspace \times \Yspace$, measurable for
$\Bor{\Xspace} \otimes \Bor{\Yspace}$, with joint law
$\jlaw \defeq \law{(X,Y)} = \pushf{(X,Y)}{\Prob}$ on the product space. We
write $\Prob_X = \pushf{X}{\Prob}$ for the input marginal.
\end{assumption}

Standardness (every uncountable such space being Borel-isomorphic to
$[0,1]$; \citealp[Thm.~A1.2]{kallenberg2002foundations}) is used only to
guarantee that the conditional laws of \cref{subsec:disint} exist as genuine
Markov kernels. No assumption is placed on $\jlaw$ beyond measurability: in
particular the empirical regime of \texttt{preprint-v1} corresponds to a
fixed but unknown $\jlaw$, and the theory makes no claim about its form.

\subsection{Representations as measurable maps}
\label{subsec:rep}

A representation is the formatting/encoding stage that turns the raw input
into the object a predictor actually consumes (the C-ladder rungs C2--C4 and
the parity control C1$^{+}$ of \texttt{preprint-v1} are concrete instances;
their categorical relationships are deferred to \cref{sec:catrep}).

\begin{definition}[Representation]
\label{def:rep}
A \emph{representation} is a measurable map $\rep : (\Xspace, \Bor{\Xspace})
\to (\Zspace_{\rep}, \Bor{\Zspace_{\rep}})$ into a standard Borel
\emph{codomain} $\Zspace_{\rep}$. Its \emph{induced sub-$\sigma$-algebra} is
\[
  \repsig \defeq \rep^{-1}\!\big(\Bor{\Zspace_{\rep}}\big)
            \;\subseteq\; \Bor{\Xspace},
\]
and its \emph{law} is the pushforward $\pushf{\rep}{\Prob_X}$ on
$\Zspace_{\rep}$. The identity map $\mathrm{id}:\Xspace\to\Xspace$ is the
maximal (lossless) representation, with $\sigalg{\mathrm{id}} = \Bor{\Xspace}$.
\end{definition}

Two facts about \cref{def:rep} are used repeatedly. First, $\repsig$ records
exactly the information $\rep$ retains: a further map $g$ factors through
$\rep$ (i.e.\ $g = h \circ \rep$ for some measurable $h$) iff $g$ is
$\repsig$-measurable (the Doob--Dynkin lemma;
\citealp[Lem.~1.13]{kallenberg2002foundations}). Second, refinement of
representations is exactly inclusion of induced $\sigma$-algebras: if $\rep'$
factors through $\rep$ then $\sigalg{\rep'} \subseteq \repsig$. This ordering
is the backbone of the data-processing inequality in \cref{sec:parity}: a
representation can only \emph{coarsen} $\Bor{\Xspace}$, never enrich it.

\begin{remark}[Lossless vs.\ lossy]
\label{rem:lossy}
A representation is \emph{lossless for $Y$} when $\repsig$ retains all of
$X$'s predictive content for $Y$ -- formally, when $\rep(X)$ is a sufficient
statistic for $Y$ (\cref{sec:parity}, made precise in T2-1). Otherwise it is
\emph{lossy}, and the rate--distortion corollary of \cref{sec:parity} bounds
the resulting risk inflation. We do not yet need the definition of
sufficiency; we only need that $\repsig$ is the object it will be stated
against.
\end{remark}

\subsection{Predictors as Markov kernels}
\label{subsec:pred}

A frozen LLM is not a deterministic function of its input: at fixed weights
it induces, through decoding, a conditional distribution over outputs. The
right abstraction is therefore a Markov kernel, not a measurable map.

\begin{definition}[Markov kernel]
\label{def:kernel}
A \emph{Markov kernel} from $(\Zspace,\Bor{\Zspace})$ to
$(\Aspace,\Bor{\Aspace})$ is a map $\kernel{K} : \Zspace \times \Bor{\Aspace}
\to [0,1]$ such that (i) $z \mapsto \kernel{K}(z, B)$ is
$\Bor{\Zspace}$-measurable for each fixed $B \in \Bor{\Aspace}$, and (ii)
$B \mapsto \kernel{K}(z, B)$ is a probability measure on $\Bor{\Aspace}$ for
each fixed $z \in \Zspace$.
\end{definition}

\begin{definition}[Predictor; frozen LLM]
\label{def:predictor}
A \emph{predictor on the representation $\rep$} is a Markov kernel
$\kernel{K} : \Zspace_{\rep} \times \Bor{\Aspace} \to [0,1]$ from the
representation codomain to an \emph{action space} $\Aspace$ (the prediction
space; $\Aspace = \Yspace$ in the simplest case). A \emph{frozen} LLM is such
a kernel held fixed: the formatting choice is absorbed into $\rep$, and
$\kernel{K}$ is the fixed map weights-plus-decoder induces from formatted
input to output distribution.
\end{definition}

A predictor on $\rep$ acts on the raw input by pre-composition with the
representation, yielding a kernel from $\Xspace$ to $\Aspace$:
\[
  (\kcomp{\kernel{K}}{\rep})(x, B) \;\defeq\; \kernel{K}\!\big(\rep(x), B\big),
  \qquad x \in \Xspace,\; B \in \Bor{\Aspace}.
\]
Because $\rep$ is measurable and $\kernel{K}$ is a kernel,
$\kcomp{\kernel{K}}{\rep}$ is again a Markov kernel
(\citealp[Ch.~2]{cinlar2011probability}); crucially it is measurable with
respect to $\repsig$, so a predictor built on $\rep$ can never depend on input
information that $\rep$ discards. This is the measure-theoretic content of
``the predictor only sees the representation,'' and it is what the
data-processing inequality will exploit.

\subsection{Loss, risk, and Bayes risk}
\label{subsec:risk}

\begin{definition}[Loss and risk]
\label{def:risk}
A \emph{loss} is a measurable map $\loss : \Yspace \times \Aspace \to
[0,\infty]$. The \emph{risk} of a predictor $\kernel{K}$ on representation
$\rep$ is the expected loss under the data law and the predictor's own
randomness,
\[
  \risk(\kernel{K}; \rep)
  \;\defeq\;
  \Expect_{(X,Y)\sim\jlaw}\;
  \int_{\Aspace} \loss\big(Y, a\big)\,
       \kernel{K}\!\big(\rep(X), \mathrm{d}a\big),
\]
the inner integral being a Lebesgue integral against the probability measure
$\kernel{K}(\rep(X),\cdot)$.
\end{definition}

\begin{definition}[Bayes risk relative to a representation]
\label{def:bayesrisk}
The \emph{Bayes risk achievable from $\rep$} is the infimum of the risk over
all predictors that act through $\rep$,
\[
  \bayesrisk(\rep)
  \;\defeq\;
  \inf_{\kernel{K}}\, \risk(\kernel{K}; \rep),
\]
the infimum ranging over all Markov kernels
$\kernel{K} : \Zspace_{\rep} \times \Bor{\Aspace} \to [0,1]$. Equivalently
(\cref{subsec:disint}) it is the risk of the best $\repsig$-measurable
decision rule. The unconstrained Bayes risk is $\bayesrisk(\mathrm{id})$.
\end{definition}

Two remarks fix the reading of \cref{def:bayesrisk}. (a)~$\bayesrisk(\rep)$ is
the \emph{information limit} of representation $\rep$: the best any predictor
could do given only what $\rep$ retains, with capacity set aside. The whole
negative half of the paper's thesis is the statement that, under parity,
$\bayesrisk(\rep)$ does not move with $\rep$ (\cref{sec:parity}); the whole
positive half is that real, capacity-bounded predictors do not attain
$\bayesrisk(\rep)$, and that this gap -- not the information limit -- is where
format can act (\cref{sec:bounded}). (b)~Because pre-composition with $\rep$
only shrinks the available $\sigma$-algebra ($\repsig \subseteq
\Bor{\Xspace}$), the map $\rep \mapsto \bayesrisk(\rep)$ is monotone under
refinement: a coarser representation has weakly larger Bayes risk. Parity is
precisely the boundary case where the inequality is tight.

\subsection{Disintegration and the conditional-expectation machinery}
\label{subsec:disint}

The arguments of \cref{sec:parity} reduce to three classical facts. We state
them in the form used downstream; proofs are standard
(\citealp[Ch.~3,~6]{folland1999real}; \citealp[Ch.~5--6]{kallenberg2002foundations}).

\begin{proposition}[Radon--Nikodym]
\label{prop:rn}
Let $\nu$ be a $\sigma$-finite measure and $\mu \ll \nu$ a measure absolutely
continuous with respect to it on a common measurable space. Then there is a
$\nu$-a.e.\ unique nonnegative measurable density $\rnd{\mu}{\nu}$ with
$\mu(A) = \int_A \rnd{\mu}{\nu}\,\mathrm{d}\nu$ for all measurable $A$.
\end{proposition}

Radon--Nikodym is what gives conditional expectation its closed form.
For an integrable random variable $W$ and a sub-$\sigma$-algebra
$\mathcal{G} \subseteq \mathcal{F}$, the conditional expectation
$\condexp{W}{\mathcal{G}}$ is the $\Prob$-a.s.\ unique
$\mathcal{G}$-measurable variable satisfying $\Expect[\,\condexp{W}{\mathcal{G}}
\,\mathbf{1}_G\,] = \Expect[W\,\mathbf{1}_G]$ for all $G \in \mathcal{G}$;
existence is \cref{prop:rn} applied to $A \mapsto \Expect[W\,\mathbf{1}_A]$ on
$\mathcal{G}$. The single property we use is the \emph{tower rule}: for nested
$\sigma$-algebras $\mathcal{G}_1 \subseteq \mathcal{G}_2$,
\begin{equation}
  \label{eq:tower}
  \condexp{\condexp{W}{\mathcal{G}_2}}{\mathcal{G}_1}
  = \condexp{W}{\mathcal{G}_1}.
\end{equation}
Applied to $\mathcal{G}_1 = \repsig \subseteq \mathcal{G}_2 = \Bor{\Xspace}$,
\cref{eq:tower} is the data-processing inequality in embryo: coarsening to
$\repsig$ is a conditional-expectation projection, and projections cannot
increase the information available to a convex risk. \cref{sec:parity} turns
this into the irrelevance theorem.

\begin{proposition}[Disintegration / regular conditional probability]
\label{prop:disint}
Under \cref{ass:sbs}, for the sub-$\sigma$-algebra $\repsig$ there exists a
Markov kernel $\post{\rep} : \Zspace_{\rep} \times \Bor{\Yspace} \to [0,1]$,
the \emph{regular conditional law of $Y$ given $\rep(X)$}, such that
\[
  \post{\rep}\!\big(\rep(X), B\big)
  = \Prob\big(Y \in B \given \repsig\big)
  \qquad \Prob\text{-a.s., for every } B \in \Bor{\Yspace}.
\]
It is $\pushf{\rep}{\Prob_X}$-a.s.\ unique. The kernel $\post{\rep}$ is the
\emph{Bayes-optimal predictor} on $\rep$: for any loss $\loss$ admitting a
measurable Bayes act, plugging $\post{\rep}$ into \cref{def:risk} attains
$\bayesrisk(\rep)$ of \cref{def:bayesrisk}.
\end{proposition}

Existence in \cref{prop:disint} is exactly where standardness of $\Yspace$ is
spent (\citealp[Thm.~6.3]{kallenberg2002foundations};
\citealp[Ch.~IV]{cinlar2011probability}); on a non-standard target the
conditional law need not regularize into a kernel. \cref{prop:disint} closes
the loop opened in \cref{subsec:pred}: the abstract ``best predictor on
$\rep$'' of \cref{def:bayesrisk} is a concrete Markov kernel, the posterior
$\post{\rep}$, and comparing $\post{\rep}$ across representations under parity
is precisely the calculation \cref{sec:parity} carries out. Everything later
in the paper is an estimate of how far a frozen LLM's kernel $\kernel{K}$
sits from $\post{\rep}$, and of how that distance moves -- or, under parity,
fails to move -- with $\rep$.
